\section{A nonempty finite positive regime}\label{sec:regime}
Both rates are positive, so squaring the condition and multiplying by the positive factor $8\alpha_N/(mk^2R^2)$ give the equivalence
\begin{equation}\label{eq:finitechoice}% SRC: FA25 N06
\sigma_{\rm ext}\le\frac\nu2
\quad\Longleftrightarrow\quad
\eta_N^2\ge
\frac{8\alpha_N}{mk^2R^2}
\left(\frac{c^2d}{m}+\frac{\kappa^2m}{\beta_R}\right).
\end{equation}
For any fixed admitted geometry and fixed finite positive $c,\kappa,\alpha_N,\beta_R,G$, the right side is finite and positive. Select $\eta_N$ as its positive square root, or any larger finite number, once before choosing the phase. Keep that value constant on every resulting solution. This gives a nonempty regime with genuinely positive propagation and receiver coupling, including the equality case.

The normalized $R^2=3$ is independent of phase, and $k,m,d$ depend only on the fixed registration and profiles. Hence both~\eqref{eq:finitechoice} and $t_*=1/\nu$ are phase independent. The theorem compares separately prepared trajectories under the same coefficients; it does not select coefficients along a trajectory. No zero-coupling limit, infinite-coupling construction or threshold optimization is needed. Appendix~\ref{sec:units} verifies that both rates have units of inverse time and the threshold for $\eta_N^2$ is dimensionless, retaining $R^2$ in every dimensional formula.
